\documentclass[12pt]{ctexart}   % 正文字体大小为12pt（可选项：10pt, 11pt, 12pt）

\usepackage{lipsum}         % 生成随机文字

% 字体定制

% 可以直接调用系统字体
% 在终端中输入命令 fc-list ，可以查看系统字体的路径、名称等信息
\usepackage{fontspec}           
\setmainfont{Times New Roman}           % 设置英文与数字的主字体
% \setsansfont{Arial}                   % 设置无衬线字体（英文与数字的粗体字体）
% \setmonofont{Courier New}             % 设置等宽字体（英文与数字的等宽（代码）字体）

% % 设置中文主字体及其粗体字体、斜体字体
% % 若调用系统中文字体，可以在终端中输入以下命令进行查找： fc-list :lang=zh-cn
\setCJKmainfont[
    BoldFont=方正小标宋简体, 
    ItalicFont=楷体
]{霞鹜文楷等宽 GB 屏幕阅读版}


\usepackage{minted}         % 代码高亮


% 页面大小与页边距设置
\usepackage{geometry}   
\geometry{a4paper, 
    left=2.54cm, right=2.54cm, top=3.18cm, bottom=3.18cm,
    headsep=1cm,    % 页眉和正文之间的距离
    footskip=1cm    % 页脚和正文之间的距离
}

% 页眉页脚设置
\usepackage{fancyhdr}
\pagestyle{fancy}
\fancyhf{}                      % 清空所有页眉和页脚区域的当前设置，避免继承默认的或其他样式残留的设置
\setlength{\headheight}{15pt}   % 设置页眉区域的高度为 15 pt
\addtolength{\topmargin}{-2pt}  % 调整页面上边距（\topmargin）的大小，将其减少 2 pt
\fancyfoot[C]{\thepage}         % 页脚居中显示页码
\fancyhead[L]{\LaTeX{}文字格式调整}      
\fancyhead[R]{You Know Who}            


% 行距
\usepackage{setspace}
\onehalfspacing

% 设置一级标题、二级标题、三级标题的格式
\ctexset{%
  section = {
    format     = \centering\bfseries\fontsize{16bp}{20bp}\selectfont,    % 格式：居中对齐、加粗、字体大小16bp、行距20bp
    aftername  = \quad, % 编号与内容的间隔
    beforeskip = 25bp,  % 段前间距
    afterskip  = 12bp,  % 段后间距
    name   = {,.},      % 编号格式
  },
  subsection = {
    format     = \bfseries\fontsize{14bp}{18bp}\selectfont,
    aftername  = \quad,
    beforeskip = 12bp,
    afterskip  = 6bp,
  },
  subsubsection = {
    format     = \bfseries\fontsize{13bp}{15bp}\selectfont,
    aftername  = \quad,
    beforeskip = 12bp,
    afterskip  = 6bp,
  },
}%


\usepackage[hidelinks]{hyperref}


\title{\bfseries \LaTeX{}文字格式调整演示\thanks{本文档受到乌有先生之乌有基金的资助}}   % 标题
\author{You Know Who, \and You Cannot-Know Who\thanks{小组长/指导教师/通讯作者：wuyou@nosuchanemail.no}}        % 作者
\date{\today}                   % 日期


\begin{document}


\maketitle
\tableofcontents

\thispagestyle{empty}   % 封面及目录页的页眉、页脚为空

\newpage



\section{正文脚注演示}

\lipsum[1]\footnote{这是一个脚注}

\section{文本加粗演示}

\textbf{\lipsum[2]}

{\bfseries \lipsum[2]}

汉字{\bfseries 加粗}演示。

\section{文本斜体演示}

\textit{\lipsum[2]}

{\itshape \lipsum[2]}

汉字{\itshape 斜体}演示。

\section{字体调整与定制}

在宏包 \verb|fontspec|、\verb|ctex| 以及 \verb|xelatex| 编译命令下，我们可以定制汉字以及数字和英文字符的正常字体、粗体、斜体、等宽字体对应的字体。
示例代码如下：
\begin{verbatim}
\usepackage{fontspec}   % 字体定制支持

% 中文字体设置，给定字体文件
\setCJKmainfont[
    Path = my_fonts/,                           % 字体文件所在目录
    BoldFont = SourceHanSerifSC-Heavy.otf,      % 粗体对应的字体文件名
    ItalicFont = ChillKai.ttf                   % 斜体对应的字体文件名
]{SourceHanSerifSC-Regular.otf}                 % 正文字体文件名

% 中文字体设置，使用系统字体
\setCJKmainfont[
    BoldFont = 方正小标宋简体, 
    ItalicFont = 楷体
]{霞鹜文楷等宽 GB 屏幕阅读版}
\end{verbatim}

% 字体下载链接：https://cloud.tsinghua.edu.cn/f/64111188e3eb4084b4f8/?dl=1

\section{代码块演示}

在宏包 \verb|minted| 的帮助下，我们可以在\LaTeX{}中插入语法高亮的代码块。此时，上述的示例代码块可呈现如下：
\begin{minted}
[
    frame=lines,
    framesep=2mm,
    baselinestretch=1.2,
    linenos
]
{latex}
\usepackage{fontspec}   % 字体定制支持

% 中文字体设置，给定字体文件
\setCJKmainfont[
    Path = my_fonts/,                           % 字体文件所在目录
    BoldFont = SourceHanSerifSC-Heavy.otf,      % 粗体对应的字体文件名
    ItalicFont = ChillKai.ttf                   % 斜体对应的字体文件名
]{SourceHanSerifSC-Regular.otf}                 % 正文字体文件名

% 中文字体设置，使用系统字体
\setCJKmainfont[
    BoldFont = 方正小标宋简体, 
    ItalicFont = 楷体
]{霞鹜文楷等宽 GB 屏幕阅读版}
\end{minted}

在本地配置时，需有Python环境（版本不低于3.8）并安装 Pygments 库：
\begin{minted}
[
    frame=lines,
    framesep=2mm,
    baselinestretch=1.2,
    linenos
]{shell}    
pip install Pygments
\end{minted}
并且在编译命令中，需增加 \verb|-shell-escape| 参数。（注意到这一行没有缩进）一个示例VSCode配置如下：
\begin{minted}
[
    frame=lines,
    framesep=2mm,
    baselinestretch=1.2,
    linenos
]{json}    
{
    "name": "xelatex",
    "command": "xelatex",
    "args": [
        "-synctex=1",
        "-shell-escape",
        "-interaction=nonstopmode",
        "-file-line-error",
        "%DOCFILE%"
    ]
},
\end{minted}

However, Overleaf can save you the trouble of installing it and having to run special commands to compile your document---on Overleaf, documents that use minted will work ``out of the box". 
(See \href{https://www.overleaf.com/learn/latex/Code_Highlighting_with_minted}{here} for detailed guides.)

请注意，英文中一对引号需表示为：\verb|``content"|，其呈现的效果是：``content"。

\section{下划线演示}

这是\underline{下划线文字}。


\section{文字大小调整演示}

那座山，正当顶上，有{\zihao{1}一块仙石}。其石有{\zihao{3}三}丈{\zihao{6}六}尺{\zihao{-5}五}寸高，有{\zihao{2}二}丈{\zihao{4}四}尺围圆\footnote{三丈六尺五寸高，按周天三百六十五度；二丈四尺围圆，按政历二十四气。}。
上有九窍八孔，按九宫八卦。四面更无树木遮阴，左右倒有芝兰相衬。


% 右对齐，左对齐对应flushleft环境
\begin{flushright}
    ——节选自《西游记》第一回
\end{flushright}

\section{列表演示}
\subsection{有序列表（enumerate环境）演示}

问：把大象装进冰箱一共需要几步？

答：三步。详细步骤如下，仅供参考：
\begin{enumerate}
    \item 把冰箱门打开
    \item 把大象装进去
    \item 把冰箱门带上
\end{enumerate}

更具体的步骤如下，仅仅仅仅供参考：
\begin{enumerate}
    \item 把冰箱门打开
    \item 把大象装进去
    \begin{enumerate}
        \item 把大象压缩
        \item 把大象移动到冰箱门前
        \item 把大象通过冰箱门放入冰箱中
    \end{enumerate}
    \item 把冰箱门带上
\end{enumerate}

\subsection{无序列表（itemize环境）演示}

\begin{itemize}
    \item Nam dui ligula, fringilla a, euismod sodales, sollicitudin vel, wisi. 
    \item Morbi auctor lorem non justo. Nam lacus libero, pretium at, lobortis vitae, ultricies et, tellus. 
    \item Donec aliquet, tortor sed accumsan bibendum, erat ligula aliquet magna, vitae ornare odio metus a mi. 
\end{itemize}


\section{\LaTeX{}中的特殊字符}

\# \$ \% \& \{ \} \_ \textbackslash

对应的\LaTeX{}代码如下：
\begin{minted}
[
    frame=lines,
    framesep=2mm,
    baselinestretch=1.2,
    linenos
]
{latex}
\# \$ \% \& \{ \} \_ \textbackslash
\end{minted}

\end{document}


